


\begin{document}
	
	% RESET alle Zähler am Anfang
	\setcounter{section}{0}
	\setcounter{subsection}{0}
	\setcounter{subsubsection}{0}
	\setcounter{paragraph}{0}
	
	% Tiefe für Nummerierung und TOC
	\setcounter{secnumdepth}{1}  % Nur Sections nummerieren
	% Part in TOC: footnotesize, bold, NO page break
	\makeatletter
	\renewcommand*\l@part[2]{%
		\ifnum \c@tocdepth >-2\relax
		\addpenalty{-\@highpenalty}%
		\addvspace{0.8em \@plus\p@}%
		{\leftskip 0em \relax
			\rightskip \@tocrmarg
			\parfillskip -\rightskip
			\parindent 0em \relax\@afterindenttrue
			\interlinepenalty\@M
			\leavevmode
			{\footnotesize\bfseries #1}\nobreak
			\leaders\hbox{$\m@th\mkern \@dotsep mu\hbox{}\mkern \@dotsep mu$}\hfill
			\nobreak\hb@xt@\@pnumwidth{\hss #2}\par}%
		\addvspace{0.2em \@plus\p@}%
		\nobreak
		\fi}
	\makeatother
	
	% Chapter in TOC: footnotesize, bold
	\renewcommand{\cftchapfont}{\footnotesize\bfseries}
	\renewcommand{\cftchappagefont}{\footnotesize\bfseries}
	\setlength{\cftbeforechapskip}{0.3em}
	
	% Only Chapters in TOC (no Sections/Subsections)
	\setcounter{tocdepth}{0}
	

	\begin{titlepage}
		\centering
		\vspace*{2cm}
	{\Huge\textbf{Fundamental Fractal-Geometric Field Theory (FFGFT) or T0 Theory – Time-Mass Duality}}\\[1cm]
	{\Large Part 7: The Lattice in Hilbert Space, Spectral Theory and Galois Structure}\\[2cm]
		{\large Johann Pascher}\\[1cm]
		{\large 2026}
		\vfill
	\end{titlepage}

	\frontmatter
	\pagestyle{fancy}
	\makeatletter
	\let\ps@plain\ps@fancy
	\let\ps@empty\ps@fancy
	\makeatother

	\mainmatter
	\pagestyle{fancy}
	\tableofcontents

	\chapter{Introduction to Volume 7}
	% =============================================================================
% INTRODUCTION TO VOLUME 7: THE LATTICE IN HILBERT SPACE, SPECTRAL THEORY AND
% GALOIS STRUCTURE (Vol. 7 of 8)
% =============================================================================

\section*{Seventh Volume of the Document Collection}

Volume~7 continues the extension of the collection begun in Volume~6 and
comprises Documents 311 to 343 from August and early September 2026. In this
period the focus of the corpus shifts from the projection into observable
space to the algebraic structure of the compact geometry itself.

\subsection*{Content and Focus}

The first part begins with the question of how a four-dimensional lattice
arrives at three spatial dimensions (Doc.~311), and leads via the closure
scale, in which $\Lambda$ appears as a derived quantity of the compact
geometry, and the thermal history on the time cycle to the lattice in Hilbert
space (Doc.~314). There the lattice is read as the index set of the Fourier
modes; geometry becomes spectrum and fibre. The investigation of the form of
$K_\text{frak}$ and the comparison with the Riemann zeta function complete
this part.

The second part brings together the lepton generations in a topological
reading, the scope of the mass derivations, the spectral theory of the FFGFT
tori with its Hilbert-space embedding, the algebraic derivation of the
$SU(3)_c$ gauge structure from $\mathbb{Z}_3$ triality, and the derivation of
the Weinberg angle at $M_Z$ with the complete renormalisation-group running.
The third part contains the comparison with the hyperbit vacuum structure
$G(6,\mathbb{Z}_3\mathbb{C})$, the FFGFT Hawking mechanism, the
self-adjointness of the fundamental operator $\hat F$, the coupling regimes of
resonance and synchronisation, and the two roads from the discrete to the
continuum.

The fourth part contains the Galois bridge. Starting from the finite field
$GF(9)$ (Doc.~336) it is shown that the Frobenius automorphism and the FFGFT
sector pairing are structurally identical; via $GF(27)^*$ there follow the
comparison of the mass ratios with Galois group orders, the Frobenius
separation of the massive and massless sectors, the neutrino mass hierarchy,
the factorisation classes in $G(6)$ with the harmonics of the primes, and the
zeta function of the $T^4/\mathbb{Z}_3$ torus.

\subsection*{Organisation}

The 33 documents are divided into four parts: four onto three, the closure
scale and the lattice in Hilbert space (Docs.~311--316); leptons, spectral
theory and gauge structure (Docs.~317--323); hyperbits, the Hawking
mechanism, coupling and the continuum (Docs.~324--335); the Galois bridge
(Docs.~336--343).

\subsection*{Reading the Documents}

Documents 311 to 314 were originally typeset as multi-chapter treatises; in
the book their chapters appear as sections of the respective document. Status
markers and the corrections register apply as in the preceding volumes; in
each case the most recent version of a result in the corpus is authoritative.

\subsection*{Status}

The documents in this volume reflect their respective versions in the corpus
as of October 2026.


	% --- part ---
	\part{Part I --- Four onto Three, the Closure Scale and the Lattice in Hilbert Space}
	\FFGFTdocDemote{Doc.~311 --- Four onto Three --- How does a 4D lattice arrive at three spatial dimensions?}{../ch/311_Vier_auf_Drei_En_ch}
	\FFGFTdocDemote{Doc.~312 --- The Closure Scale: $\Lambda$ as a Derived Quantity of the Compact Geometry}{../ch/312_Abschlussskala_En_ch}
	\FFGFTdoc{Appendix A to Doc.~312 --- Altermagnetism as a Laboratory Analogy}{../ch/312_Anhang_Altermagnet_En_ch}
	\FFGFTdocDemote{Doc.~313 --- No Beginning: The Thermal History on the Time Cycle}{../ch/313_Kein_Anfang_En_ch}
	\FFGFTdocDemote{Doc.~314 --- The Lattice in Hilbert Space: Geometry, Spectrum and Fibre}{../ch/314_Gitter_im_Hilbertraum_En_ch}
	\FFGFTdoc{Doc.~315 --- The Form of $K_\text{frak}$: Additive or Multiplicative?}{../ch/315_Kfrak_Form_En_ch}
	\FFGFTdoc{Doc.~316 --- FFGFT and the Riemann Zeta Function}{../ch/316_Riemann_Zeta_En_ch}
	% --- part ---
	\part{Part II --- Leptons, Spectral Theory and Gauge Structure}
	\FFGFTdoc{Doc.~317 --- Topological Origin of Lepton Generations: KSAU and FFGFT}{../ch/317_KSAU_FFGFT_Leptonen_En_ch}
	\FFGFTdoc{Doc.~318 --- Scope of Mass Derivations in FFGFT: Rest Masses, Hadron Dynamics, and Open Bridges}{../ch/318_Masseableitung_Geltungsbereich_En_ch}
	\FFGFTdoc{Doc.~319 --- The Proton as a Vibrating Torus: Geometric Foundations of Mass Emergence in FFGFT}{../ch/319_Proton_Torus_En_ch}
	\FFGFTdoc{Doc.~320 --- Detailed Field-Geometric Spectral Theory}{../ch/320_Spektraltheorie_En_ch}
	\FFGFTdoc{Doc.~321 --- Algebraic Derivation of the $SU(3)_c$ Gauge Structure}{../ch/321_SU3_Z3_Emergenz_En_ch}
	\FFGFTdoc{Doc.~322 --- Spectral Theory and Hilbert-Space Embedding of the FFGFT Tori}{../ch/322_Spektraltheorie_Hilbert_En_ch}
	\FFGFTdoc{Doc.~323 --- Derivation of the Weinberg Angle at $M_Z$ from FFGFT Geometry}{../ch/323_Weinberg_Winkel_RGE_En_ch}
	% --- part ---
	\part{Part III --- Hyperbits, the Hawking Mechanism, Coupling and the Continuum}
	\FFGFTdoc{Doc.~324 --- $G(6,\mathbb{Z}_3\mathbb{C})$ Vacuum Structure and FFGFT}{../ch/324_G6_Z3C_FFGFT_Vergleich_En_ch}
	\FFGFTdoc{Doc.~325 --- The FFGFT Hawking Mechanism}{../ch/325_Hawking_FFGFT_En_ch}
	\FFGFTdoc{Doc.~326 --- Black Holes from Hyperbits and from FFGFT}{../ch/326_Matzke_FFGFT_Vergleich_En_ch}
	\FFGFTdoc{Doc.~327 --- Self-Adjointness of the Fundamental Fractal Operator $\hat{F}$}{../ch/327_Selbstadjungiertheit_F_En_ch}
	\FFGFTdoc{Doc.~328 --- Coupling Regimes, Resonance, and Synchronisation}{../ch/328_Kopplungsregime_Resonanz_En_ch}
	\FFGFTdoc{Doc.~329 --- Is the Stability Threshold $n_{\mathrm{thresh}}$ Universal?}{../ch/329_nthresh_Skalenanalyse_En_ch}
	\FFGFTdoc{Doc.~330 --- Three Operations from $T^4$ to Observable Space}{../ch/330_Operationen_T4_Abgrenzung_En_ch}
	\FFGFTdoc{Doc.~332 --- Two Roads from the Discrete to the Continuum}{../ch/332_Kontinuumskonstruktionen_En_ch}
	\FFGFTdoc{Doc.~333 --- $K_{\text{frak}}$, Unrolling and the $\tilde{T}\cdot m$ Duality}{../ch/333_Kfrak_Dualitaet_En_ch}
	\FFGFTdoc{Doc.~334 --- Superposition Without Time}{../ch/334_Superposition_ohne_Zeit_En_ch}
	\FFGFTdocDemote{Doc.~335 --- Mass-Bound Time Scales and Interferometry}{../ch/335_Frequenzskalen_En_ch}
	% --- part ---
	\part{Part IV --- The Galois Bridge: GF(9), GF(27) and the Zeta Function of the Torus}
	\FFGFTdoc{Doc.~336 --- The Algebraic Bridge: FFGFT and $\mathbb{Z}_3\mathbb{C}$}{../ch/336_GF9_FFGFT_Bruecke_En_ch}
	\FFGFTdoc{Doc.~337 --- Time Emergence: FFGFT and GALG Compared}{../ch/337_Zeit_Emergenz_FFGFT_GALG_En_ch}
	\FFGFTdoc{Doc.~338 --- Mass Calculations and Galois Comparison}{../ch/338_Galois_Massen_FFGFT_En_ch}
	\FFGFTdoc{Doc.~339 --- Frobenius Separation of Massive and Massless Sectors in $GF(27)^*$}{../ch/339_Frobenius_Trennung_En_ch}
	\FFGFTdoc{Doc.~340 --- Neutrino Mass Hierarchy from the $GF(27)^*$ Frobenius Structure}{../ch/340_Neutrino_Galois_En_ch}
	\FFGFTdoc{Doc.~341 --- GF$(27)$ in GALG and FFGFT: Why G$(3)$ Contains No GF$(27)$ but G$(6)$ Does}{../ch/341_GF27_GALG_FFGFT_En_ch}
	\FFGFTdoc{Doc.~342 --- Factorisation Classes in G$(6)$ and the Harmonics of the Primes}{../ch/342_Faktorisierung_Harmonik_En_ch}
	\FFGFTdoc{Doc.~343 --- The Zeta Function of the T$^4$/Z$_3$ Torus: Orbifold Projection and Dirichlet L-Functions}{../ch/343_Zeta_Galois_FFGFT_En_ch}
\end{document}
